%------------------------------------------------------------------------------%
\documentclass[fleqn,utf8,aspectratio=169,12pt]{beamer}

\usepackage{calculo_varias_variaveis-1}

\title{Fórmulas de De Moivre}

%------------------------------------------------------------------------------%
\begin{document}

\addtitle

%------------------------------------------------------------------------------%
\section{Potencias de Números Complexos}

%------------------------------------------------------------------------------%
\begin{proofframe}{Justificativa}
  \begin{align*}
    \onslide<+->{z^n}
    \onslide<+->{& = zz\cdots z \\[5mm]}
    \onslide<+->{& = \rho\rho\cdots\rho}
    \onslide<+->{    \big[\cos(\varphi+\varphi+\cdots+\varphi)
                       + i\sin(\varphi+\varphi+\cdots+\varphi)
                     \big] \\[5mm]}
    \onslide<+->{& = \rho^n\big[\cos(n\varphi) + i\sin(n\varphi)\big]}
  \end{align*}
\end{proofframe}

%------------------------------------------------------------------------------%
\begin{frame}{Primeira Fórmula de De Moivre}
  Se
  \[
    z = \rho \big[\cos(\varphi) + i\sin(\varphi)\big]
  \]
  e $n$ um número natural

  \pause
  \vspace{\baselineskip}
  Primeira Fórmula de De Moivre
  \[
    z^n = \rho^n \big[\cos(n\varphi) + i\sin(n\varphi)\big]
  \]
\end{frame}

%------------------------------------------------------------------------------%
\include{L-05-Formulas_Moivre-ex1}
\include{L-05-Formulas_Moivre-ex2}

%------------------------------------------------------------------------------%
\section{Raízes de Números Complexos}

%------------------------------------------------------------------------------%
\begin{proofframe}{Justificativa}
  Se
  \(
    z = \rho \big[\cos(\varphi) + i\sin(\varphi)\big]
  \)
  e $n>1$ é um número natural

  \pause
  \vspace{\baselineskip}
  A raiz $n$-ésima de $z$ é um número
  \[
    u = r \big[\cos(\alpha) + i\sin(\alpha)\big]
  \]
  \pause
  tal que
  \[
    u^n = z
  \]
  \[
    \pause
    r^n \big[\cos(n\alpha) + i\sin(n\alpha)\big]
    = \rho \big[\cos(\varphi) + i\sin(\varphi)\big]
  \]
\end{proofframe}

%------------------------------------------------------------------------------%
\begin{proofframe}{Justificativa}
\begin{columns}
\column{0.5\linewidth}
  Encontrar $r$ e $\alpha$
  \pause
   tais que
  \begin{gather*}
    r^n = \rho \\
    \cos(n\alpha) = \cos(\varphi) \\
    \sin(n\alpha) = \sin(\varphi)
  \end{gather*}
  \pause
  ou seja
  \begin{gather*}
    r = \sqrt[n]{\rho} \\
    n\alpha = \varphi + 2k\pi
  \end{gather*}
  \column{0.5\linewidth}
  \pause
  \[
    \alpha = \frac{\varphi}{n} + 2\pi\frac{k}{n}
  \]
  \pause
  apenas
  \[
    k = 0, 1, 2, \dots, n-1
  \]
  fornecem soluções distintas
\end{columns}
\end{proofframe}

%------------------------------------------------------------------------------%
\begin{frame}{Segunda Formula de De Moire}
  Se
  \(
    z = \rho \big[\cos(\varphi) + i\sin(\varphi)\big]
  \)
  e $n>1$ é um número natural

  \vspace{\baselineskip}
  As raízes $n$-ésimas de $z$ são
  \[
    u_k = \sqrt[n]{\rho}
    \left[
         \cos\left(\frac{\varphi}{n}+2\pi\frac{k}{n}\right)
      + i\sin\left(\frac{\varphi}{n}+2\pi\frac{k}{n}\right)
    \right]
  \]
  com \(k=0, 1, 2, \dots, n-1\)
\end{frame}

%------------------------------------------------------------------------------%
\section{Exemplos}

\include{L-05-Formulas_Moivre-ex3}
\include{L-05-Formulas_Moivre-ex4}
\include{L-05-Formulas_Moivre-ex5}
\include{L-05-Formulas_Moivre-ex6}
\include{L-05-Formulas_Moivre-ex7}

%------------------------------------------------------------------------------%
\section{Lista Mínima}

%------------------------------------------------------------------------------%
\begin{frame}{Lista Mínima}
  Atenção: A prova é baseada no livro, não nas apresentações
\end{frame}

%------------------------------------------------------------------------------%
\end{document}
%------------------------------------------------------------------------------%
